\documentclass[11pt, a4paper]{article}

% ── Packages ──────────────────────────────────────────────────────
\usepackage[utf8]{inputenc}
\usepackage[T1]{fontenc}
\usepackage{lmodern}
\usepackage[margin=1in]{geometry}
\usepackage{amsmath, amssymb, amsthm}
\usepackage{mathtools}
\usepackage{booktabs}
\usepackage{array}
\usepackage{hyperref}
\usepackage{enumitem}
\usepackage{xcolor}
\usepackage{microtype}
\usepackage{tikz-cd}

% ── Theorem environments ──────────────────────────────────────────
\theoremstyle{plain}
\newtheorem{theorem}{Theorem}[section]
\newtheorem{lemma}[theorem]{Lemma}
\newtheorem{proposition}[theorem]{Proposition}

\theoremstyle{definition}
\newtheorem{definition}[theorem]{Definition}
\newtheorem{remark}[theorem]{Remark}

% ── Custom commands ───────────────────────────────────────────────
\newcommand{\Agda}{\textsc{Agda}}
\newcommand{\HoTT}{\textsc{HoTT}}
\newcommand{\RT}{Ryu--Takayanagi}
\newcommand{\AdS}{AdS}
\newcommand{\dS}{dS}
\newcommand{\CFT}{CFT}
\newcommand{\ua}{\mathsf{ua}}
\newcommand{\transport}{\mathsf{transport}}
\newcommand{\refl}{\mathsf{refl}}
\newcommand{\funExt}{\mathsf{funExt}}
\newcommand{\isoToEquiv}{\mathsf{isoToEquiv}}
\newcommand{\Qgeq}{\mathbb{Q}_{\geq 0}}
\newcommand{\Qten}{\mathbb{Q}_{10}}
\newcommand{\NN}{\mathbb{N}}
\newcommand{\ZZ}{\mathbb{Z}}
\newcommand{\Type}{\mathsf{Type}_0}
\newcommand{\isSet}{\mathsf{isSet}}
\newcommand{\isProp}{\mathsf{isProp}}
\newcommand{\isContr}{\mathsf{isContr}}
\newcommand{\Scut}{S_{\mathrm{cut}}}
\newcommand{\Lmin}{L_{\mathrm{min}}}
\newcommand{\Obs}{\mathsf{Obs}}
\newcommand{\BW}{\mathsf{BridgeWitness}}
\newcommand{\PD}{\mathsf{PatchData}}
\newcommand{\ORP}{\mathsf{OrbitReducedPatch}}

% ── Title ─────────────────────────────────────────────────────────
\title{%
  \textbf{Univalence Gravity} \\[6pt]
  \large Project Capstone \& Formalization Summary \\[4pt]
  \normalsize A Constructive Formalization of Discrete
  Entanglement--Geometry Duality in Cubical Agda}

\author{Sven Bichtemann}

\date{April 2026 \\ \smallskip
  \small Compiler: Agda 2.8.0 \quad $\mid$ \quad
  Library: \texttt{agda/cubical} (April 2026) \quad $\mid$ \quad
  Paradigm: Homotopy Type Theory}

% ══════════════════════════════════════════════════════════════════
\begin{document}
\maketitle

% ──────────────────────────────────────────────────────────────────
\begin{abstract}
\noindent
The \emph{Univalence Gravity} project is a constructive
formalization effort in Cubical \Agda{} that investigates whether
finite, discrete models of the holographic correspondence from
quantum gravity can be stated and verified as exact type
equivalences in Homotopy Type Theory (\HoTT{}).  Inspired by the
Anti-de~Sitter / Conformal Field Theory (\AdS/\CFT{}) duality and the
\RT{} formula---which equates boundary entanglement entropy with
bulk minimal-surface area---the project builds a family of
combinatorial toy models (HaPPY-code patches on hyperbolic tilings
and honeycombs), extracts boundary and bulk observable packages as
dependent record types, and constructs machine-checked type
equivalences $\Obs_\partial \simeq \Obs_{\mathrm{bulk}}$ whose
$\ua$-paths yield computable transport functions.  Three core
theorems have been mechanized: (1)~a discrete Gauss--Bonnet theorem
for combinatorial curvature on the $\{5,4\}$ tiling, verified by
\refl{} on a class-weighted integer sum; (2)~structural properties
(subadditivity of boundary entropy, monotonicity of bulk chain
length) proven by exhaustive case analysis with witnesses sealed
behind Agda's \texttt{abstract} barrier; and (3)~exact enriched
observable-package equivalences with verified $\transport$ along
$\ua$-paths, demonstrated across twelve patch instances spanning
1D~trees, 2D pentagonal tilings, and 3D cubic honeycombs.
A generic bridge theorem (\texttt{GenericBridge.agda}),
parameterized by an abstract \PD{} interface, factors the proof
into a single reusable kernel instantiated at each geometric scale
by a Python oracle, enabling the formalization to scale to patches
with thousands of tiles without additional hand-written proof.
The project further demonstrates a discrete Wick rotation
(curvature-sign flip between $\{5,4\}$ and $\{5,3\}$ tilings
sharing the same flow graph), a resolution tower with monotone
min-cut spectrum growth, and parameterized step-invariance of the
holographic correspondence under local bond-weight perturbations.
The backend mathematical foundation is now complete, paving the
way for Haskell extraction of the verified translator functions and
interactive WebGL visualization.
\end{abstract}

\tableofcontents

% ══════════════════════════════════════════════════════════════════
\section{Introduction and Motivation}
\label{sec:intro}

The AdS/CFT correspondence, proposed by Maldacena in 1997, posits
a deep equivalence between quantum gravity in an
Anti-de~Sitter bulk spacetime and a conformal field theory on its
boundary.  The \RT{} formula sharpens this into a precise equation:
the entanglement entropy $S(A)$ of a boundary region~$A$ equals the
area of the minimal bulk surface homologous to~$A$, measured in
Planck units.  These correspondences have never been stated as
equivalences between precisely defined mathematical objects in any
formal system, nor has the ``translation'' between dual descriptions
ever been extracted as a verified program.

This project asks a narrow, well-defined question:

\begin{quote}
\emph{Can a finite, discrete toy model of a holographic
correspondence be formalized in Cubical \Agda{} such that boundary
and bulk observable packages are connected by an exact type
equivalence, with the Univalence Axiom providing computational
transport between them?}
\end{quote}

The contribution is positioned as a result in \textbf{proof
engineering and constructive formalization}, not as a new physical
theory.  All formalized objects are finite and combinatorial; no
continuum limits, smooth manifolds, or path integrals appear.  The
vehicle is Homotopy Type Theory, where the Univalence Axiom
identifies equivalences with paths between types, and Cubical
\Agda{} gives $\transport$ along such paths genuine computational
content.

\subsection{Scope and Non-Goals}

The following are explicitly out of scope:
\begin{itemize}[nosep]
  \item No continuum limits or smooth geometry.
  \item No full tensor-network quantum mechanics (Hilbert spaces,
    complex amplitudes).
  \item No dependence on AI for any formal result.
  \item No claim about physical spacetime---results concern
    mathematical toy models.
  \item No cohesive HoTT implementation (discussed as a future
    conceptual horizon only).
\end{itemize}

% ══════════════════════════════════════════════════════════════════
\section{Mathematical Architecture}
\label{sec:architecture}

\subsection{The Common Source / Observable Package Pattern}

Rather than attempting to map raw structural types to each other,
the formalization relies on a \textbf{Common Source Specification}
$c : C$ which encodes the tiling, triangulation, and weight
assignments.  From this source, two distinct views are extracted:

\begin{enumerate}[nosep]
  \item \textbf{Boundary view} $\pi_\partial(c)$: a finite discrete
    entanglement network paired with a min-cut entropy functional
    $\Scut : \mathsf{Region} \to \Qgeq$.
  \item \textbf{Bulk view} $\pi_{\mathrm{bulk}}(c)$: a finite
    simplicial or cell complex paired with a minimal separating
    chain/surface functional $\Lmin : \mathsf{Region} \to \Qgeq$.
\end{enumerate}

These functionals are packaged into \emph{observable packages}---record types parameterized by a region index type.  Pointwise agreement
$\Scut(r) \equiv \Lmin(r)$ for every region~$r$, assembled via
$\funExt$ into a function-level path, is the discrete \RT{}
correspondence.

\subsection{The Enriched Equivalence Construction}

The specification-agreement enrichment pairs each observable
function with a proof that it matches a reference specification:
\[
  \mathsf{EnrichedBdy} \;=\; \Sigma_{f : R \to \Qgeq}\,(f \equiv S_\partial),
  \qquad
  \mathsf{EnrichedBulk} \;=\; \Sigma_{f : R \to \Qgeq}\,(f \equiv L_B).
\]
These are reversed singleton types (contractible), but genuinely
different types in the universe because $S_\partial$ and $L_B$ are
definitionally distinct functions defined by separate pattern
matches.  An explicit $\mathsf{Iso}$ is built by appending the
\RT{} path $\mathsf{obs\text{-}path} : S_\partial \equiv L_B$ to
the agreement witness; round-trip proofs close because $R \to \Qgeq$
is a set ($h$-level~2), making specification-agreement fields
propositional. Promotion via $\isoToEquiv$ and application of
$\ua$ yields the Univalence path; $\transport$ reduces via
$\mathsf{ua}\beta$ to the concrete forward map.

\subsection{The Scalar Representation}

Two scalar types serve distinct roles:
\begin{itemize}[nosep]
  \item $\Qgeq = \NN$ (\texttt{Util/Scalars.agda}): nonnegative
    scalars for min-cut and chain-length values.  Judgmental
    computation of $\NN$ addition enables all pointwise agreement
    and subadditivity proofs to discharge by $\refl$.
  \item $\Qten = \ZZ$ (\texttt{Util/Rationals.agda}): signed
    rationals in tenths for combinatorial curvature.  The integer
    $n$ represents $n/10$; all curvature values for $\{5,4\}$ and
    $\{5,3\}$ tilings have denominators dividing~10.
\end{itemize}

Constants are defined \emph{once} in utility modules and imported
everywhere, guaranteeing identical normal forms---the critical
property enabling $\refl$-based proofs.

% ══════════════════════════════════════════════════════════════════
\section{Core Theorems Mechanized}
\label{sec:theorems}

Three definitive theorems have been fully mechanized, each
progressing through numerical prototyping (Python oracles),
Agda code generation, and hand-written Cubical \Agda{} proof
assembly.

% ──────────────────────────────────────────────────────────────────
\subsection{Theorem 1: Discrete Gauss--Bonnet}
\label{sec:thm1}

\begin{theorem}[Discrete Gauss--Bonnet]
For the 11-tile filled patch~$K$ of the $\{5,4\}$ hyperbolic
pentagonal tiling, with combinatorial curvature
$\kappa : \mathsf{Vertex} \to \Qten$ defined by
\[
  \kappa(v) \;=\; 1 - \frac{\deg_E(v)}{2}
  + \sum_{f \ni v} \frac{1}{\mathrm{sides}(f)},
\]
the total curvature equals the Euler characteristic:
\[
  \sum_{v \in K} \kappa(v) \;=\; \chi(K) \;=\; 1.
\]
\end{theorem}

\noindent
\textbf{Proof architecture.}
The 30~vertices of the polygon complex are classified into 5
classes by combinatorial neighborhood (\texttt{Bulk/PatchComplex.agda}).
Curvature is constant within each class
(\texttt{Bulk/Curvature.agda}), so Gauss--Bonnet reduces to a
class-weighted sum over $\ZZ$:
\[
  5 \cdot (-2) + 5 \cdot (-1) + 5 \cdot 2 + 5 \cdot (-1) + 10 \cdot 2
  \;=\; 10 \;=\; \mathsf{one}_{10}.
\]
The proof is $\refl$ because $\ZZ$ addition computes by structural
recursion on closed terms (\texttt{Bulk/GaussBonnet.agda}).  A
\texttt{GaussBonnetWitness} record packages all data and the
proof as a single inspectable milestone artifact.

\medskip\noindent
\textbf{De~Sitter analogue.}
A parallel Gauss--Bonnet theorem is proven for the $\{5,3\}$
spherical tiling (regular dodecahedron star patch), where interior
curvature is \emph{positive} ($\kappa = +1/10$).  Both theorems
target $\chi = 1$ but with different interior/boundary
decompositions:
\[
  \{5,4\}\text{ (AdS)}: \; (-10) + 20 = 10,
  \qquad
  \{5,3\}\text{ (dS)}: \; (+5) + 5 = 10.
\]

% ──────────────────────────────────────────────────────────────────
\subsection{Theorem 2: Structural Properties---Subadditivity and Monotonicity}
\label{sec:thm2}

\begin{theorem}[Boundary Subadditivity \& Bulk Monotonicity]
For the 6-tile star patch of the $\{5,4\}$ HaPPY code:
\begin{enumerate}[nosep,label=(\alph*)]
  \item \textup{(Subadditivity)} For every representable union
    $r_1 \cup r_2 = r_3$ among 20 boundary regions:
    $\Scut(r_3) \leq \Scut(r_1) + \Scut(r_2)$.
  \item \textup{(Monotonicity)} For every subregion inclusion
    $r_1 \subseteq r_2$ among 10 representative regions:
    $\Lmin(r_1) \leq \Lmin(r_2)$.
\end{enumerate}
\end{theorem}

\noindent
\textbf{Proof.}
The ordering $m \leq n$ is encoded as
$\Sigma_{k:\NN}\,(k + m \equiv n)$, where $k$ is a concrete
natural number and the path component is $\refl$ (since $\NN$
addition computes judgmentally on closed terms).  Subadditivity
involves 30~union triples discharged by four distinct witnesses:
$(0,\refl)$, $(1,\refl)$, $(2,\refl)$, $(3,\refl)$.
Monotonicity involves 10~subregion inclusions, all witnessed by
$(1,\refl)$.  Propositionality of $\leq$ is established by
$\NN$ right-cancellation (\texttt{isProp$\leq_\QQ$}).

The 11-tile filled patch's 360 subadditivity triples are verified
by a Python-generated module
(\texttt{Boundary/FilledSubadditivity.agda}) sealed behind
\texttt{abstract}, preventing downstream RAM cascades.

% ──────────────────────────────────────────────────────────────────
\subsection{Theorem 3: The Holographic Bridge / Observable-Package Equivalence}
\label{sec:thm3}

\begin{theorem}[Holographic Bridge]
For the 6-tile star patch of the $\{5,4\}$ HaPPY code, the
enriched boundary observable package $\mathsf{FullBdy}$ (carrying
a subadditivity witness) and the enriched bulk observable package
$\mathsf{FullBulk}$ (carrying a monotonicity witness) are exactly
equivalent as types:
\[
  \mathsf{FullBdy} \;\simeq\; \mathsf{FullBulk}.
\]
Transport along the resulting Univalence path carries the canonical
boundary instance to the canonical bulk instance:
\[
  \transport\,(\ua\,\mathsf{full\text{-}equiv})\;
  \mathsf{full\text{-}bdy} \;\equiv\; \mathsf{full\text{-}bulk}.
\]
The structural property is converted: a subadditivity witness in
the input becomes a monotonicity witness in the output.
\end{theorem}

\noindent
\textbf{Instantiation breadth.}
This theorem has been instantiated across \textbf{twelve} distinct
patches, as summarized in Table~\ref{tab:instances}.

\begin{table}[ht]
\centering
\small
\begin{tabular}{@{}llrrrl@{}}
\toprule
\textbf{Instance} & \textbf{Geom.} & \textbf{Regions} &
\textbf{Orbits} & \textbf{Max $S$} & \textbf{Strategy} \\
\midrule
Tree pilot       & 1D tree   &    8 & --- &  2 & flat refl \\
Star (6-tile)    & 2D $\{5,4\}$ &  10 & --- &  2 & flat refl \\
Filled (11-tile) & 2D $\{5,4\}$ &  90 & --- &  4 & Python gen \\
Honeycomb (BFS)  & 3D $\{4,3,5\}$ & 26 & --- &  1 & Python gen \\
Dense-50         & 3D $\{4,3,5\}$ & 139 & --- & 7 & Python gen \\
Dense-100        & 3D $\{4,3,5\}$ & 717 &  8 &  8 & orbit red. \\
Dense-200        & 3D $\{4,3,5\}$ & 1246 & 9 &  9 & orbit red. \\
$\{5,4\}$ depth~2 & 2D $\{5,4\}$ & 15 & 2 &  2 & orbit red. \\
$\{5,4\}$ depth~3 & 2D $\{5,4\}$ & 40 & 2 &  2 & orbit red. \\
$\{5,4\}$ depth~4 & 2D $\{5,4\}$ & 105 & 2 & 2 & orbit red. \\
$\{5,4\}$ depth~5 & 2D $\{5,4\}$ & 275 & 2 & 2 & orbit red. \\
$\{5,4\}$ depth~7 & 2D $\{5,4\}$ & 1885 & 2 & 2 & orbit red. \\
\bottomrule
\end{tabular}
\caption{All mechanized bridge instances.  ``Orbit red.''
  denotes the orbit reduction strategy where proofs are on the
  small orbit type and lifted to the full region type via a
  classification function.}
\label{tab:instances}
\end{table}

% ══════════════════════════════════════════════════════════════════
\section{Proof Engineering Breakthroughs}
\label{sec:engineering}

The formalization confronted several fundamental performance and
scalability barriers inherent to dependently typed proof assistants.
This section documents the obstacles encountered and the techniques
developed to overcome them.

\subsection{The AST/RAM Wall}
\label{sec:ram}

Agda's type-checker normalizes terms by recursively unpacking every
constructor into abstract syntax trees (ASTs) in memory,
structurally verifying that both sides of an equality reduce to the
same normal form.  For 360~subadditivity cases on the 11-tile
patch, each involving several levels of constructor unfolding, the
memory cost is enormous.  This is not a bug but a deliberate
architectural trade-off rooted in the \textbf{De~Bruijn Criterion}:
the trusted computing base of a proof assistant must be as small
and simple as possible.

\subsection{Python Oracles (Approach~C)}

External Python scripts enumerate all combinatorial cases, compute
the correct witnesses, and output complete Agda modules containing
explicit $(k,\refl)$ proofs.  The Python script is the
``oracle'' that \emph{finds} proofs; Agda is the ``checker'' that
\emph{verifies} each line individually.  Thirteen Python scripts
(\texttt{01}--\texttt{13} in \texttt{sim/prototyping/}) serve as
oracles for increasingly complex patches, from the 6-tile star
through 3046-tile $\{5,4\}$ depth-7 layers.

This pattern is aligned with how the largest verified proofs in
computer science were achieved: external computation to find proofs,
a simple kernel to check them (cf.\ the Four Color Theorem in Coq,
the Kepler Conjecture in HOL~Light).

\subsection{The \texttt{abstract} Barrier (Approach~A)}

Wrapping a proof block in \texttt{abstract} tells Agda: ``once
verified, seal it and never unfold it again.''  This prevents
downstream modules from re-normalizing the 360-case proof when
they use the subadditivity lemma, halting the RAM cascade
(cf.\ Agda GitHub Issue~\#4573).  Since $\leq_\QQ$ is propositional
(by \texttt{isProp$\leq_\QQ$}), sealing the proof does not damage
the Univalence bridge: propositional fields are transported
trivially, and the enriched equivalence depends only on
$\mathsf{obs\text{-}path}$ (which is \emph{not} abstract and
retains full computational content).

\subsection{Orbit / Symmetry Reduction}
\label{sec:orbit}

The Dense-100 patch has 717 boundary regions but only 8 distinct
min-cut values.  The \textbf{orbit reduction strategy} classifies
regions into orbits by min-cut value, defines observable lookups on
the small orbit type, and lifts the agreement proof via a single
line:
\[
  \mathsf{d100\text{-}pointwise}\; r \;=\;
  \mathsf{d100\text{-}pointwise\text{-}rep}\;(\mathsf{classify100}\; r).
\]
This reduces proof obligations from 717 to 8.  The classification
function (717 clauses) is traversed only during concrete evaluation,
never during proof normalization.  The Dense-200 instance
(1246 regions $\to$ 9 orbits), and $\{5,4\}$ depth-7 instance
(1885 regions $\to$ 2 orbits), confirm that the pattern scales
to arbitrary sizes: adding tiles grows only the classify function,
not the proof obligations.

\subsection{Schematic Bridge Factorization}
\label{sec:generic}

The decisive architectural innovation of the project is the
\textbf{Generic Bridge Theorem}
(\texttt{Bridge/GenericBridge.agda}), a parameterized module
proving that \emph{any} \PD{} record admits the full enriched
equivalence, Univalence path, and verified transport.

\begin{definition}[\PD{} Interface]
A $\PD$ record captures:
\begin{itemize}[nosep]
  \item $\mathsf{RegionTy} : \Type$ (the region type),
  \item $S_\partial, L_B : \mathsf{RegionTy} \to \Qgeq$ (the two
    observable functions),
  \item $\mathsf{obs\text{-}path} : S_\partial \equiv L_B$ (the
    discrete \RT{} correspondence).
\end{itemize}
\end{definition}

\begin{definition}[\ORP{} Interface]
An $\ORP$ refines $\PD$ with orbit data:
$\mathsf{OrbitTy}$, $\mathsf{classify}$, $S_{\mathrm{rep}}$,
$L_{\mathrm{rep}}$, and $\mathsf{rep\text{-}agree}$.  The
conversion $\mathsf{orbit\text{-}to\text{-}patch}$ extracts the
$\PD$ automatically, constructing $\mathsf{obs\text{-}path}$ by
$\funExt$ of the lifted orbit-level agreement.
\end{definition}

The composition
$\mathsf{orbit\text{-}bridge\text{-}witness} :
\ORP \to \BW$
produces the full proof-carrying bridge witness in one function
call.  This module is \textbf{written once, proven once, and never
modified}.  Every bridge instance in Table~\ref{tab:instances} is a
special case, as verified retroactively by
\texttt{Bridge/GenericValidation.agda} (all six pre-existing bridges
validated).

The \texttt{SchematicTower} module
(\texttt{Bridge/SchematicTower.agda}) assembles verified
\texttt{TowerLevel} records (each containing an $\ORP$ + max
min-cut + $\BW$), connected by \texttt{LayerStep}
monotonicity witnesses.  The full $\{5,4\}$ tower spans depths
2--7 (21 to 3046 tiles), and the Dense tower spans 100--200 cells
with growing min-cut spectrum $7 \to 8 \to 9$.

% ══════════════════════════════════════════════════════════════════
\section{Additional Formalized Results}
\label{sec:additional}

Beyond the three core theorems, several additional results have
been mechanized.

\subsection{Discrete Wick Rotation}

The \texttt{WickRotationWitness} record
(\texttt{Bridge/WickRotation.agda}) packages:
\begin{itemize}[nosep]
  \item the shared holographic bridge (Theorem~3, curvature-agnostic),
  \item the AdS Gauss--Bonnet witness ($\kappa_{\mathrm{int}} = -1/5$,
    hyperbolic),
  \item the dS Gauss--Bonnet witness ($\kappa_{\mathrm{int}} = +1/10$,
    spherical),
  \item a coherence proof that both target the same
    $\chi_{10} = \mathsf{pos}\;10$.
\end{itemize}
The ``discrete Wick rotation'' is a parameter change in the tiling
type ($q = 4 \to q = 3$) that flips the curvature sign while
preserving the flow graph.  No complex numbers, no analytic
continuation, no Lorentzian geometry.

\subsection{Coarse-Graining and the Discrete Area Law}

The \texttt{CoarseGrainedRT} type
(\texttt{Bridge/CoarseGrain.agda}) packages the orbit reduction
as a thermodynamic coarse-graining: an observer with only
$\lceil \log_2 8 \rceil = 3$ bits of memory reads the correct
entropy value for each orbit class.  The discrete isoperimetric
inequality $\Scut(A) \leq \mathrm{area}(A)$ is verified on all
717~regions of the Dense-100 patch and all 1246~regions of
Dense-200 (\texttt{Boundary/Dense100AreaLaw.agda},
\texttt{Boundary/Dense200AreaLaw.agda}), with each case
sealed in \texttt{abstract}.

\subsection{Parameterized Step-Invariance and Dynamics}

The step-invariance theorem
(\texttt{Bridge/StarStepInvariance.agda}) proves that the discrete
\RT{} correspondence $S_{\mathrm{param}}\,w \equiv
L_{\mathrm{param}}\,w$ is preserved under single-bond weight
perturbations for \emph{arbitrary} $w : \mathsf{Bond} \to \Qgeq$.
The ``while(true) loop'' (\texttt{Bridge/StarDynamicsLoop.agda})
extends this by structural induction on lists of perturbation
steps, establishing that arbitrarily long sequences of local
modifications maintain the holographic correspondence.  The
enriched version
(\texttt{Bridge/EnrichedStarStepInvariance.agda}) further proves
that the full $\mathsf{FullBdy}(w) \simeq \mathsf{FullBulk}(w)$
equivalence (with subadditivity $\leftrightarrow$ monotonicity
conversion) holds for any weight function.

\subsection{Raw Structural Equivalence}

For both the 6-tile star and the 11-tile filled patch,
raw structural equivalences are constructed between the
15-dimensional bulk bond-weight type and the 90-dimensional
boundary min-cut type (\texttt{Bridge/StarRawEquiv.agda},
\texttt{Bridge/FilledRawEquiv.agda}).  The forward map is the
discrete \RT{} computation; the backward map uses a peeling
strategy (G-singleton extraction + truncated subtraction) to
recover bond weights from min-cut values.

\subsection{Resolution Tower and Convergence Certificate}

A 3-level \texttt{ConvergenceCertificate3L}
(\texttt{Bridge/ResolutionTower.agda}) records:
\begin{center}
\begin{tabular}{@{}lrrrl@{}}
\toprule
\textbf{Level} & \textbf{Regions} & \textbf{Orbits} &
\textbf{Max $S$} & \textbf{Monotone} \\
\midrule
Dense-50  & 139  & ---  & 7 & --- \\
Dense-100 & 717  &   8 & 8 & $(1,\refl)$ \\
Dense-200 & 1246 &   9 & 9 & $(1,\refl)$ \\
\bottomrule
\end{tabular}
\end{center}
The min-cut spectrum grows monotonically ($7 \leq 8 \leq 9$),
and the area-law bound holds at every level.

% ══════════════════════════════════════════════════════════════════
\section{Repository Architecture}
\label{sec:repo}

The source code is organized into four layers:

\begin{description}[style=nextline,nosep]
  \item[\texttt{Util/}] Scalar types ($\Qgeq = \NN$,
    $\Qten = \ZZ$) and arithmetic lemmas.
  \item[\texttt{Common/}] Observable package record, per-instance
    specification types (region types, orbit types, classification
    functions).
  \item[\texttt{Boundary/} and \texttt{Bulk/}] Per-instance
    observable lookups, curvature, Gauss--Bonnet, subadditivity,
    monotonicity, area-law bounds.
  \item[\texttt{Bridge/}] Observable agreement
    ($\mathsf{obs\text{-}path}$), enriched equivalences, Univalence
    paths, transport verification, the generic bridge, validation,
    schematic tower, Wick rotation, coarse-graining, dynamics.
\end{description}

The dependency DAG after the generic factorization is:
\[
  \underbrace{\text{Spec} \to \text{Cut/Chain} \to \text{Obs}}_{\text{per-instance (oracle-generated)}}
  \;\longrightarrow\;
  \underbrace{\text{GenericBridge}}_{\text{proven once}}
  \;\longrightarrow\;
  \underbrace{\text{SchematicTower}}_{\text{all instances registered}}
\]
The hand-written per-instance Equiv modules
(\texttt{EnrichedStarEquiv}, \texttt{FilledEquiv}, etc.)\ are
architecturally superseded by the generic bridge but remain valid
as independently verified artifacts.

A companion simulation layer (\texttt{sim/prototyping/}) contains
13 Python scripts that serve as oracles for graph construction
(Coxeter reflections for $\{5,4\}$ and $\{4,3,5\}$), growth
strategies (BFS, greedy dense, geodesic tube), min-cut computation
(max-flow via NetworkX), orbit classification, area-law calculation,
and Agda module emission.  The environment is pinned by
\texttt{shell.nix} (NixOS 24.05, Python 3.12, NetworkX 3.6.1,
NumPy 2.4.4).

% ══════════════════════════════════════════════════════════════════
\section{Quantitative Summary}
\label{sec:quantitative}

\begin{itemize}[nosep]
  \item \textbf{Hand-written Agda modules:} $\sim$30 (Util, Common,
    core Boundary/Bulk/Bridge modules).
  \item \textbf{Oracle-generated Agda modules:} $\sim$50
    (Spec/Cut/Chain/Obs/AreaLaw for 12 patch instances).
  \item \textbf{Total Agda lines:} $\sim$25{,}000+
    (including generated code).
  \item \textbf{Python oracle scripts:} 13, totaling
    $\sim$6{,}000 lines.
  \item \textbf{Largest single type:} \texttt{D200Region}
    (1246 constructors, parsed and type-checked by Agda without
    $\mathsf{Fin}\;N$ encoding).
  \item \textbf{Largest abstract proof:} 1246 cases in
    \texttt{Boundary/Dense200AreaLaw.agda}.
  \item \textbf{Deepest $\{5,4\}$ layer:} depth 7, 3046 tiles,
    1885 regions, 2 orbits.
  \item \textbf{Highest min-cut:} $S = 9$ (Dense-200, 3D
    $\{4,3,5\}$ honeycomb).
\end{itemize}

% ══════════════════════════════════════════════════════════════════
\section{Relationship to Prior Work}
\label{sec:related}

The HaPPY code (Pastawski et al., 2015) provides the physical toy
model; this project abstracts to its combinatorial content.  Regge
calculus (1961) inspires the discrete curvature formulation.  The
$\{5,4\}$ and $\{4,3,5\}$ Coxeter geometries are standard
(Coxeter, \emph{Regular Polytopes}, 1973).  The proof engineering
techniques---external oracles, \texttt{abstract} barriers, orbit
reduction---draw on community knowledge around Agda's performance
characteristics (Agda Issue~\#4573; Vezzosi et al., FSCD 2024 on
``Automating Boundary Filling in Cubical Agda'').

To our knowledge, no prior work has:
\begin{enumerate}[nosep]
  \item constructed a machine-checked type equivalence between
    boundary and bulk observable packages for a holographic toy
    model in any proof assistant;
  \item demonstrated that the discrete \RT{} correspondence is
    curvature-agnostic (the ``discrete Wick rotation'') in a
    formally verified setting;
  \item scaled a Cubical \Agda{} formalization of a combinatorial
    physics correspondence to 3D hyperbolic cell complexes with
    over 1000 boundary regions.
\end{enumerate}

% ══════════════════════════════════════════════════════════════════
\section{Conclusion and Next Steps}
\label{sec:conclusion}

The backend mathematical foundation of the Univalence Gravity
project is complete.  The three core theorems are mechanized, the
generic bridge subsumes all per-instance constructions, and the
schematic tower demonstrates scalability to exponentially growing
hyperbolic patches.  The factorization into
$\PD / \ORP / \mathsf{GenericEnriched}$ cleanly separates geometry
(handled by the Python oracle) from proof (handled by the generic
theorem, proven once).

\subsection{Immediate Next Steps}

\begin{enumerate}[nosep]
  \item \textbf{Architectural tightening.} Factor
    \texttt{BridgeWitness} into its own module; consolidate
    \texttt{ResolutionTower} into \texttt{SchematicTower}; rewire
    \texttt{WickRotation} to reference the generic bridge witness.
  \item \textbf{Extraction entry point.} Write a \texttt{Main.agda}
    module importing all theorem statements and bridge witnesses,
    serving as the compilation target for
    \texttt{agda --compile --ghc}.
  \item \textbf{Haskell backend.} Extract the verified
    $\mathsf{equivFun}$ from generic bridge witnesses as
    standalone executable translator functions.
  \item \textbf{WebGL frontend.} Build a Three.js interactive
    application hooking into the extracted Haskell binary, allowing
    users to toggle between shallow and deep topologies and
    visualize the holographic correspondence interactively.
\end{enumerate}

\subsection{Deferred Research Directions}

Several directions have been evaluated and deliberately set aside
with detailed feasibility analyses preserved in
\texttt{docs/10-frontier.md}:

\begin{itemize}[nosep]
  \item \textbf{Full quantum mechanics:} formalizing Hilbert spaces
    and tensor contraction in Agda would require constructive
    complex analysis---an active, largely unsolved area.
  \item \textbf{Coinductive continuum limit:} defining an infinite
    stream of verified patches requires a generic patch constructor
    (equivalent to solving the N-layer problem inductively in Agda),
    plus \texttt{--guardedness}.
  \item \textbf{Causal structure:} defining a discrete light cone
    on the cell complex is a research-level contribution to causal
    set theory.
\end{itemize}

\bigskip\noindent
The project demonstrates that formal type-theoretic tools---in
particular Univalence and computational transport in Cubical
\Agda{}---can be brought to bear on a correspondence previously
discussed only informally in the theoretical physics literature.
Even the partial results achieved constitute novel contributions to
the formalization community: a verified discrete Gauss--Bonnet
theorem, a machine-checked min-cut entropy functional, and a
computable transport term that converts boundary quantum-informational
observables into bulk geometric observables across twelve distinct
geometric instantiations spanning three spatial dimensions.

% ══════════════════════════════════════════════════════════════════
\section*{Acknowledgments}

The \texttt{agda/cubical} library developers and the Agda core team
provided the foundational infrastructure.  The HaPPY code by
Pastawski, Yoshida, Harlow, and Preskill provided the physical toy
model.  NetworkX and NumPy powered the Python oracle layer.

% ══════════════════════════════════════════════════════════════════

\begin{thebibliography}{99}

\bibitem{maldacena97}
J.~Maldacena.
\newblock The large $N$ limit of superconformal field theories and
  supergravity.
\newblock \emph{Adv.\ Theor.\ Math.\ Phys.}, 2:231--252, 1998.
\newblock arXiv:hep-th/9711200.

\bibitem{rt06}
S.~Ryu and T.~Takayanagi.
\newblock Holographic derivation of entanglement entropy from
  AdS/CFT.
\newblock \emph{Phys.\ Rev.\ Lett.}, 96:181602, 2006.
\newblock arXiv:hep-th/0603001.

\bibitem{happy15}
F.~Pastawski, B.~Yoshida, D.~Harlow, and J.~Preskill.
\newblock Holographic quantum error-correcting codes: toy models
  for the bulk/boundary correspondence.
\newblock \emph{JHEP}, 2015(6):149, 2015.
\newblock arXiv:1503.06237.

\bibitem{cchm18}
C.~Cohen, T.~Coquand, S.~Huber, and A.~M\"ortberg.
\newblock Cubical type theory: a constructive interpretation of
  the univalence axiom.
\newblock \emph{J.\ Autom.\ Reason.}, 60(2):159--210, 2018.

\bibitem{vezzosi19}
A.~Vezzosi, A.~M\"ortberg, and A.~Abel.
\newblock Cubical Agda: a dependently typed programming language
  with univalence and higher inductive types.
\newblock \emph{ICFP}, 2019.

\bibitem{regge61}
T.~Regge.
\newblock General relativity without coordinates.
\newblock \emph{Nuovo Cim.}, 19:558--571, 1961.

\bibitem{coxeter73}
H.\,S.\,M.~Coxeter.
\newblock \emph{Regular Polytopes}.
\newblock Dover, 3rd edition, 1973.

\bibitem{gonthier05}
G.~Gonthier.
\newblock A computer-checked proof of the four-colour theorem.
\newblock Unpublished manuscript, Microsoft Research, 2005.

\bibitem{hottbook}
The Univalent Foundations Program.
\newblock \emph{Homotopy Type Theory: Univalent Foundations of
  Mathematics}.
\newblock Institute for Advanced Study, 2013.

\end{thebibliography}

\end{document}